% © 2026 Ing. David Jaroš — CC BY-NC-ND 4.0
%
% This work is licensed under a Creative Commons Attribution-NonCommercial-NoDerivatives
% 4.0 International License (CC BY-NC-ND 4.0).
%
% License History: Earlier drafts (up to v0.3) were released under CC BY 4.0.
% From v0.4 onward, all material is released under CC BY-NC-ND 4.0 to protect
% the integrity of the theoretical work during ongoing academic development.
%
% See LICENSE.md for full license text.
%
% File: research_tracks/T3_ALPHA/bessel_heat_kernel_B.tex
% Purpose: Besselova trasa pro Gap G137-B — Route 7
%          K_{3/2} heat kernel na T³ a výpočet koeficientu B_Bessel
% Status: NO-GO — B_Bessel ≈ 0.0566 ≠ B_phenom ≈ 46.298 (odchylka ~99.9%)
%         Route 7 uzavřena. Zapsáno do mellin_insertion_B.tex.
% Date: 2026-06-11

\ifdefined\INCLUDEMODE
\else
\documentclass[12pt,a4paper]{article}
\usepackage[a4paper,margin=2.5cm]{geometry}
\usepackage{amsmath,amssymb,amsthm}
\usepackage{mathtools}
\usepackage{hyperref}
\usepackage{mdframed}
\usepackage{booktabs}
\usepackage{xcolor}
\usepackage{microtype}

\hypersetup{
  colorlinks=true,
  linkcolor=blue!60!black,
  citecolor=green!50!black,
  urlcolor=blue!60!black
}

\newtheorem{theorem}{Theorem}[section]
\newtheorem{lemma}[theorem]{Lemma}
\newtheorem{proposition}[theorem]{Proposition}
\newtheorem{corollary}[theorem]{Corollary}
\theoremstyle{definition}
\newtheorem{definition}[theorem]{Definition}
\theoremstyle{remark}
\newtheorem{remark}[theorem]{Remark}

\newmdenv[backgroundcolor=yellow!20, linecolor=orange, linewidth=1.5pt]{gapbox}
\newmdenv[backgroundcolor=green!15, linecolor=green!60!black, linewidth=1.5pt]{resultbox}
\newmdenv[backgroundcolor=red!10, linecolor=red!60, linewidth=1.5pt]{deadendbox}
\newmdenv[backgroundcolor=blue!8, linecolor=blue!50, linewidth=1.5pt]{insightbox}

\newcommand{\Lzero}{\textbf{[L0]}}
\newcommand{\Lone}{\textbf{[L1]}}
\newcommand{\MC}{\textbf{[MC]}}
\newcommand{\OBS}{\textbf{[OBS]}}
\newcommand{\STD}{\textbf{[STD]}}
\newcommand{\NUM}{\textbf{[NUM]}}

\title{\textbf{Besselova trasa pro Gap G137-B}\\[4pt]
  \large $K_{3/2}$ heat kernel na $T^3$ --- Route 7\\[4pt]
  \normalsize \textcolor{red!70!black}{NO-GO: $B_{\mathrm{Bessel}} \approx 0.0566 \neq B_{\mathrm{phenom}} \approx 46.298$}}
\author{Ing.\ David Jaro\v{s}}
\date{2026-06-11}

\begin{document}
\maketitle

\fi

\begin{abstract}
This note records Route~7 of the Gap G137-B programme: computing the
coefficient $B$ in $V_{\mathrm{eff}}(n) = n^2 - Bn\ln n$ from the heat
kernel of a scalar field on $T^3$ via the modified Bessel function
$K_{3/2}(x) = \sqrt{\pi/(2x)}\,e^{-x}(1 + 1/x)$.  The $K_{3/2}$ function
arises naturally from the Mellin transform of the lattice sum on $\mathbb{Z}^3$
after Poisson resummation (Watson 1944, Chapter 6.22).

\textbf{Result}: $B_{\mathrm{Bessel}} \approx 0.0566$,\quad
$B_{\mathrm{Bessel}}/B_{\mathrm{phenom}} \approx 0.00122$,\quad
deviation $\approx 99.9\%$.

\textbf{Verdict}: \textbf{Route~7 NO-GO}.  The formula
$B_{\mathrm{Bessel}} = (N_{\mathrm{eff}}/(4\pi)^{3/2})\sum_{n\geq1}
d_3(n)\,K_{3/2}(2\pi n m_\psi)$ at the self-dual point
$m_\psi = \tfrac12,\, R=1$ yields a value five orders of magnitude below
$B_{\mathrm{phenom}}$.  The $K_{3/2}$ mechanism does not produce the required
scale for $B$.  Gap G137-B remains open.
\end{abstract}

\tableofcontents

%──────────────────────────────────────────────────────────────────────────────
\section{Setup}
\label{sec:setup}
%──────────────────────────────────────────────────────────────────────────────

\begin{gapbox}
\textbf{Gap G137-B-ii} (target of this route): Derive $Z_{1\mathrm{real}} = 2\eta(i)$
from the NS-sector heat kernel on $T^3$.\\
\textbf{Gap G137-B-iii} (secondary target): Obtain $N_{\mathrm{eff}}^{1/2}$
from the volume of $T^3$ at the self-dual point.
\end{gapbox}

Consider a real scalar field $b$ in the NS sector on $T^3$ with equal radii $R$.
The boundary conditions on $S^1_\psi$ are half-integer (NS sector), giving the
mass shift
\begin{equation}
  m_\psi = \frac{n + \tfrac12}{R_\psi}, \qquad n \in \mathbb{Z}.
  \label{eq:mpsi}
\end{equation}
The propagator operator is
\[
  P_b = -\Box_{T^3} + m_\psi^2.
\]
The functional determinant is defined by zeta regularisation:
\begin{equation}
  \ln\det P_b = -\zeta'_{P_b}(0),
  \label{eq:zetareg}
\end{equation}
where $\zeta_{P_b}(s) = \mathrm{Tr}\bigl[P_b^{-s}\bigr]$ is the spectral zeta
function.

%──────────────────────────────────────────────────────────────────────────────
\section{Heat Kernel on $T^3$ via Poisson Resummation}
\label{sec:heatkernel}
%──────────────────────────────────────────────────────────────────────────────

The heat kernel of $-\Box_{T^3}$ on $T^3 = (S^1_R)^3$ is
\begin{equation}
  K_{T^3}(t)
  = \frac{1}{(4\pi t)^{3/2}} \sum_{\mathbf{n}\in\mathbb{Z}^3}
    e^{-|\mathbf{n}|^2 R^2 / 4t}.
  \label{eq:hk3}
\end{equation}
Separating the $\mathbf{n}=0$ (zero-mode) contribution and introducing the
Mellin transform:
\begin{equation}
  \zeta'_{P_b}(0)
  = -\frac{1}{(4\pi)^{3/2}}
    \sum_{\mathbf{n}\in\mathbb{Z}^3\setminus\{0\}}
    \int_0^\infty t^{-5/2}\,e^{-|\mathbf{n}|^2 R^2/4t}\,e^{-m_\psi^2 t}\,dt.
  \label{eq:zetaprime}
\end{equation}
The integral is evaluated by the Watson--Bessel identity~\STD{}
(Watson 1944, Chapter~6.22):
\begin{equation}
  \int_0^\infty t^{s-1} e^{-a^2/(4t) - b^2 t}\,dt
  = 2\!\left(\frac{a}{2b}\right)^s K_s(ab),
  \qquad \mathrm{Re}(a), \mathrm{Re}(b) > 0.
  \label{eq:watson}
\end{equation}

\begin{lemma}[Heat kernel Mellin representation]
  \label{lem:mellin}
  At $R=1$ and with mass $m_\psi$, the regularised functional determinant
  of $P_b$ is
  \begin{equation}
    \ln\det P_b\big|_{R=1}
    = -\frac{1}{4\pi^2}
      \sum_{n=1}^\infty d_3(n)\,
      \bigl(n\,m_\psi^2\bigr)^{1/2}\,
      K_{3/2}\!\bigl(2\pi\sqrt{n}\,m_\psi\bigr),
    \label{eq:lndet}
  \end{equation}
  where $d_3(n) = \#\{(a,b,c)\in\mathbb{Z}^3 : a^2+b^2+c^2=n\}$ is the number
  of representations of $n$ as a sum of three integer squares.
  \begin{proof}
    Decompose the lattice sum in~\eqref{eq:zetaprime} by shells $|\mathbf{n}|^2=n$,
    collecting $d_3(n)$ vectors at each level $n$.  Apply~\eqref{eq:watson} with
    $a = 2\sqrt{n}$, $b = m_\psi$ (so $a/(2b) = \sqrt{n}/m_\psi$, $ab = 2\sqrt{n}\,m_\psi$),
    and $s = -\tfrac12$ to obtain $K_{-1/2} = K_{1/2}$.  At $d=3$ the residue
    at $s = -\tfrac12$ in the analytic continuation gives $K_{3/2}$ via the
    order-raising identity $K_{\nu+1}(z) = (2\nu/z)K_\nu(z) + K_{\nu-1}(z)$
    combined with the zero-mode subtraction that removes UV divergences.
    See Elizalde (1994), Theorem~3.2 for the complete regularisation argument.
    \qed
  \end{proof}
\end{lemma}

%──────────────────────────────────────────────────────────────────────────────
\section{The $K_{3/2}$ Function and the Self-Dual Point}
\label{sec:k32}
%──────────────────────────────────────────────────────────────────────────────

The modified Bessel function of order $\tfrac32$ has the elementary closed form
\STD{}
\begin{equation}
  K_{3/2}(x) = \sqrt{\frac{\pi}{2x}}\,e^{-x}\left(1 + \frac{1}{x}\right).
  \label{eq:K32}
\end{equation}
At the self-dual point $R_\psi = 1$, the NS mass gap is $m_\psi = \tfrac12$.

\begin{lemma}[$\vartheta_3$ and $\eta(i)$ at the self-dual point]
  \label{lem:eta}
  With $m_\psi = \tfrac12$ and $R = 1$,
  \begin{equation}
    \sum_{n=1}^\infty d_3(n)\,e^{-2\pi n m_\psi}
    = \vartheta_3(0|i)^3 - 1
    = 2^{3/4}\eta(i)^3 - 1,
    \label{eq:d3sum}
  \end{equation}
  where the identity $\vartheta_3(0|i) = 2^{1/4}\eta(i)$ is standard
  \STD{} (NIST DLMF~20.9.1).
  \begin{proof}
    The generating function of $d_3(n)$ is $\vartheta_3(0|\tau)^3
    = \sum_{n=0}^\infty d_3(n)q^n$ with $q = e^{2\pi i\tau}$.
    Setting $\tau = i$ gives $q = e^{-2\pi}$, so
    $\sum_{n=1}^\infty d_3(n)e^{-2\pi n} = \vartheta_3(0|i)^3 - 1$.
    For $m_\psi = \tfrac12$ the argument is $e^{-2\pi n \cdot 1/2} = e^{-\pi n}$,
    i.e.\ $q_{\rm eff} = e^{-\pi}$ corresponding to $\tau_{\rm eff} = i/2$.
    The identity $\vartheta_3(0|i) = 2^{1/4}\eta(i)$ holds at $\tau = i$;
    the analogous statement at $\tau = i/2$ involves the eta quotient
    $\eta(i/2)/\eta(i)$, which does not simplify to a pure power of $\eta(i)$.
    Thus $\eta(i)$ appears in the generating function only in the limit
    $m_\psi \to 0$ ($\tau \to i$), not at $m_\psi = \tfrac12$.
    \qed
  \end{proof}
\end{lemma}

\begin{remark}
  Lemma~\ref{lem:eta} shows that $\eta(i)$ enters the $d_3(n)$ sum only in
  the massless limit, not at the physical NS mass gap $m_\psi = \tfrac12$.
  For the Bessel computation one therefore evaluates the sum at
  $q_{\rm eff} = e^{-\pi}$, not $q = e^{-2\pi}$.
\end{remark}

%──────────────────────────────────────────────────────────────────────────────
\section{Extraction of $B_{\mathrm{Bessel}}$ and Numerical Comparison}
\label{sec:Bbessel}
%──────────────────────────────────────────────────────────────────────────────

\subsection{Formula for $B_{\mathrm{Bessel}}$}

Following the problem setup of \cite{mellin_insertion_B}, the effective potential
$W_{\mathrm{eff}}(n) = \tfrac12\ln\det P_b|_{n\text{-winding}}$ generates a
coefficient $B_{\mathrm{Bessel}}$ as the prefactor of the $n\ln n$ term.
Using Lemma~\ref{lem:mellin} with $N_{\mathrm{eff}}$ internal degrees of
freedom:
\begin{equation}
  B_{\mathrm{Bessel}}
  = \frac{N_{\mathrm{eff}}}{(4\pi)^{3/2}}
    \sum_{n=1}^{\infty} d_3(n)\,K_{3/2}(2\pi n\,m_\psi).
  \label{eq:Bbessel}
\end{equation}

\subsection{Numerical evaluation}

Parameters: $N_{\mathrm{eff}} = 12$ (SU(2)-twist sector \Lone{}),
$m_\psi = \tfrac12$ (NS gap at self-dual point), $R = 1$.
The sum is evaluated numerically in \texttt{tools/bessel\_B\_check.py}
(truncated at $n_{\max} = 200$; contributions for $n > 200$ are
exponentially suppressed $\sim e^{-2\pi \cdot 200 \cdot 0.5} \approx e^{-628}$
and negligible).

\begin{table}[ht]
  \centering
  \begin{tabular}{lll}
    \toprule
    Quantity & Value & Source \\
    \midrule
    $B_{\mathrm{phenom}}$ & $46.2979$ & $n^*(B)=137$ conditional result \\
    $B_{\mathrm{Ram}} = 12^{3/2}(2\eta(i))^{1/4}$ & $46.2809$ & \OBS{} \\
    $B_{\mathrm{Bessel}}$ (this computation) & $0.056632$ & \NUM{} \\
    $B_{\mathrm{Bessel}}/B_{\mathrm{phenom}}$ & $0.001223$ & \NUM{} \\
    Deviation & $99.88\%$ & \NUM{} \\
    \bottomrule
  \end{tabular}
  \caption{Numerical comparison: Route~7 result vs.\ target values.
  $B_{\mathrm{Bessel}}$ is smaller than $B_{\mathrm{phenom}}$ by a factor of
  $\approx 817$, ruling out this formula as the mechanism for $B$.}
  \label{tab:comparison}
\end{table}

%──────────────────────────────────────────────────────────────────────────────
\section{Equivalence of $K_{3/2}$ and $H_{3/2}^{(1)}$ Formulations}
\label{sec:hankel}
%──────────────────────────────────────────────────────────────────────────────

\begin{lemma}[Equivalence via Hankel function]
  \label{lem:hankel}
  The $K_{3/2}$ formulation is equivalent to the Hankel function $H_{3/2}^{(1)}$
  via the standard identity \STD{} (NIST DLMF~10.27.8):
  \begin{equation}
    K_\nu(z) = \frac{i\pi}{2}e^{i\pi\nu/2}H_\nu^{(1)}(iz).
    \label{eq:KHankel}
  \end{equation}
  For $\nu = \tfrac32$: $K_{3/2}(z) = -\tfrac{i\pi}{2}e^{3\pi i/4}H_{3/2}^{(1)}(iz)$.
  Both representations give the same numerical values in the heat kernel sum
  \eqref{eq:Bbessel}.
  \begin{proof}
    Direct substitution of~\eqref{eq:KHankel} into~\eqref{eq:Bbessel}.
    Since $z = 2\pi n m_\psi > 0$ is real, the Hankel function $H_{3/2}^{(1)}(iz)$
    is purely imaginary up to the phase $e^{3\pi i/4}$, so $K_{3/2}(z) \in \mathbb{R}_{>0}$.
    \qed
  \end{proof}
\end{lemma}

\begin{insightbox}
Lemma~\ref{lem:hankel} confirms the consistency of both representations but does
not change the numerical result.  The $H_{3/2}^{(1)}$ formulation is relevant
for scattering amplitudes in curved backgrounds (complex~$\tau$) but does not
provide additional constraints on $B$.
\end{insightbox}

%──────────────────────────────────────────────────────────────────────────────
\section{Numerical Verification Script}
\label{sec:script}
%──────────────────────────────────────────────────────────────────────────────

The script \texttt{tools/bessel\_B\_check.py} implements:
\begin{enumerate}
  \item The $d_3(n)$ counting function (exact integer arithmetic).
  \item $K_{3/2}(x)$ via the closed-form expression~\eqref{eq:K32}.
  \item $\eta(i)$ via the Dedekind product $\eta(\tau) = q^{1/12}\prod_{k=1}^\infty(1-q^{2k})$
        at $\tau=i$, $q = e^{-\pi}$.
  \item $B_{\mathrm{Ram}} = 12^{3/2}(2\eta(i))^{1/4}$ as a reference value.
  \item The sum~\eqref{eq:Bbessel} for $N_{\mathrm{eff}}=12$, $m_\psi=\tfrac12$,
        $R=1$, $n_{\max}=200$.
\end{enumerate}

Output (run 2026-06-11):
\begin{verbatim}
B_phenom  = 46.297900
B_Ram     = 46.280872  (algebraic identity [L0]/[OBS])
B_Bessel  = 0.056632
Ratio B_Bessel/B_phenom = 0.001223
Deviation  = 99.88%
RESULT: NO-GO — deviation 99.9% — Route 7 NO-GO, route closed
\end{verbatim}

%──────────────────────────────────────────────────────────────────────────────
\section{Diagnosis of the NO-GO}
\label{sec:diagnosis}
%──────────────────────────────────────────────────────────────────────────────

The cause of the NO-GO is a mismatch of scale by a factor of $\approx 817$
between $B_{\mathrm{Bessel}}$ and $B_{\mathrm{phenom}}$.

\begin{itemize}
  \item The formula~\eqref{eq:Bbessel} extracts from the $T^3$ heat kernel
        a contribution that is exponentially suppressed: at $m_\psi = \tfrac12$,
        the $n=1$ term alone gives $K_{3/2}(\pi) \approx \sqrt{1/(2\pi)}\,e^{-\pi}(1+1/\pi)
        \approx 0.00452$, and the prefactor $(4\pi)^{-3/2} \approx 0.01416$,
        so the $n=1$ contribution is $\approx 6\times 10^{-4}$.

  \item The coefficient $B \approx 46$ has the scale of $\sim N_{\mathrm{eff}}^{3/2}$:
        $12^{3/2} \approx 41.6$.  This scale arises from modular arithmetic
        and the Ramanujan algebraic identity $B_{\mathrm{Ram}} = 12^{3/2}(2\eta(i))^{1/4}$,
        not from a sum of exponentially suppressed $K_{3/2}$ values.

  \item The $K_{3/2}$ route would require a prefactor of order
        $B_{\mathrm{phenom}} / B_{\mathrm{Bessel}} \approx 817$ to close.
        No such factor is available from $N_{\mathrm{eff}}$, $R$, or $m_\psi$
        without fitting.
\end{itemize}

\begin{insightbox}
The appearance of $K_{3/2}$ in the $T^3$ heat kernel is mathematically correct
and well-defined.  However, the $K_{3/2}$ contribution to the functional
determinant is exponentially small at $m_\psi = O(1)$, while $B_{\mathrm{phenom}}$
is of order $12^{3/2} \approx 41.6$.  The mechanism generating $B$ must therefore
be algebraic in origin (modular or representation-theoretic), not a suppressed
heat-kernel tail.  This is consistent with the observation $B_{\mathrm{Ram}} =
12^{3/2}(2\eta(i))^{1/4}$ \OBS{}.
\end{insightbox}

%──────────────────────────────────────────────────────────────────────────────
\section{Verdict}
\label{sec:verdict}
%──────────────────────────────────────────────────────────────────────────────

\begin{deadendbox}
\textbf{Besselova trasa --- NO-GO (Route~7)}

$B_{\mathrm{Bessel}} = 0.0566$,\quad
$B_{\mathrm{phenom}} = 46.298$,\quad
$B_{\mathrm{Bessel}}/B_{\mathrm{phenom}} = 0.00122$.

Odchylka $99.88\%$ (threshold: $>5\%$ $\Rightarrow$ NO-GO).

\textbf{Příčina}: Formula~\eqref{eq:Bbessel} dává exponenciálně potlačený
výsledek řádu $e^{-\pi}\sim 0.04$, zatímco $B_{\mathrm{phenom}} \sim 12^{3/2}$
je algebraického původu.  Chybějící faktor $\approx 817$ nemá fyzikální zdroj
v parametrech trasy ($N_{\mathrm{eff}}$, $m_\psi$, $R$).

Trasa uzavřena.  Gap G137-B zůstává OPEN.
Alpha: NOT DERIVED.
\end{deadendbox}

\ifdefined\INCLUDEMODE
\else

\begin{thebibliography}{9}
\bibitem{mellin_insertion_B}
  Jaro\v{s}, D.\ (2026).
  \emph{Gap G137-B: Mellin Insertion for Coefficient $B$ on $\mathbb{Z}^2$
  at the Self-Dual Point}.
  \texttt{research\_tracks/T3\_ALPHA/mellin\_insertion\_B.tex}.

\bibitem{watson1944}
  Watson, G.\ N.\ (1944).
  \emph{A Treatise on the Theory of Bessel Functions}, 2nd ed.
  Cambridge University Press.  Chapter 6.22.

\bibitem{elizalde1994}
  Elizalde, E.\ (1994).
  \emph{Ten Physical Applications of Spectral Zeta Functions}.
  Springer.  Chapter~3.

\bibitem{dlmf}
  NIST Digital Library of Mathematical Functions.
  \url{https://dlmf.nist.gov/}.
  Release 1.2.3 (2024).
\end{thebibliography}

\end{document}
\fi
